\documentclass[11pt,a4paper]{article}

\usepackage{amsmath,amssymb,amsthm}
\usepackage{hyperref}
\usepackage[margin=1in]{geometry}
\usepackage{xcolor}
\usepackage{booktabs}
\usepackage{array}
\usepackage{float}
\usepackage{mathtools}

\newtheorem{theorem}{Theorem}[section]
\newtheorem{conjecture}[theorem]{Conjecture}
\newtheorem{corollary}[theorem]{Corollary}
\newtheorem{lemma}[theorem]{Lemma}
\newtheorem{proposition}[theorem]{Proposition}
\newtheorem{definition}[theorem]{Definition}
\theoremstyle{remark}
\newtheorem{remark}[theorem]{Remark}

\newcommand{\PT}{\mathbb{PT}}
\newcommand{\C}{\mathbb{C}}
\newcommand{\Z}{\mathbb{Z}}
\newcommand{\F}{\mathbb{F}}
\newcommand{\cO}{\mathcal{O}}
\newcommand{\cC}{\mathcal{C}}
\newcommand{\PP}{\mathcal{P}}
\newcommand{\Wop}[1]{W(#1)}


\title{The Weyl Product Identity and the Algebraic Origin\\
       of the Kochen--Specker Obstruction\\
       {\large (Paper XVI)}}

\author{Zhou-Li Chen \\ co-nlang Research}
\date{June 2026}

\begin{document}

\maketitle

\begin{abstract}
We resolve the context-level $\beta$ identification problem
(Paper~XV Open Problem~1) by showing that $\beta(C)/2 \bmod 2$ is
the accumulated phase of the Weyl operator product
$W_C := W(v_1)W(v_2)W(v_3)W(v_4)$,
where $W(v)$ is the Weyl (Pauli-string) operator for ray $v \in \mathbb{F}_2^6$.
Our main result is:

\begin{center}
\textbf{Weyl Product Identity:}~
$W_C \;=\; s(C) \cdot I_8$
\end{center}

for every context $C = \{v_1, v_2, v_3, v_4\}$,
where $s(C) = (-1)^{\beta(C)/2}$ is the context sign from Paper~XI.
The proof uses three ingredients:
(1)~the \emph{Pauli commutation relation}
$W(v)W(w) = (-1)^{\omega(v,w)} W(w)W(v)$, which for Lagrangian pairs
($\omega=0$) gives exact commutativity;
(2)~the \emph{Fano zero-sum property}
$v_1 \oplus v_2 \oplus v_3 \oplus v_4 = 0$ (an algebraic identity in the
Fano-plane structure of each Lagrangian), which forces
$W_C \in \{\pm I_8\}$;
and (3)~a \emph{direct phase computation} on $|0\rangle$ showing that
the scalar equals $(-1)^{\beta(C)/2}$ where
$\beta(C) = \sum_{j<k} \omega_{\mathrm{int}}(v_j, v_k)$.
The Kochen--Specker obstruction follows immediately:
$\prod_{C \in \mathbf{P}} W_C = -I_8$
for all 12{,}096 Mermin pentagrams,
recovering $\beta_{\mathrm{sum}} \equiv 2 \pmod 4$
as a Weyl-algebra identity.
This resolves the open question from Papers~XI--XV:
the mod~4 structure is not a lift of a mod-2 phenomenon but
is encoded in the integer symplectic form $\omega_{\mathrm{int}}$
via the $\beta$-formula, whose sign is realised as the
Pauli-string product $W_C$.
\end{abstract}


\section{Introduction}

Papers~XI through~XV established the parity theorem for Mermin pentagrams
and systematically explored its origins
\cite{PaperXI,PaperXII,PaperXIII,PaperXIV,PaperXV}.
For each context $C = \{v_1,v_2,v_3,v_4\}$, the context sign
\[
  s(C) = (-1)^{\beta(C)/2}, \qquad
  \beta(C) = \sum_{j<k} \omega_{\mathrm{int}}(v_j, v_k),
\]
satisfies $\prod_{i=1}^5 s(C_i) = -1$ for all 12{,}096 Mermin pentagrams
\cite{PaperXI}.
Paper~XV established that $\beta(C)$ is a \emph{context-level} invariant,
depending on the specific four rays of $C$ and not just the
Lagrangian $L_C$, and identified the problem of finding a
group-theoretic object encoding $\beta(C)$ at context level as the
primary open problem for further work \cite{PaperXV}.

This paper solves that problem.
The key observation is that the Weyl (Pauli-string) operators on $\mathbb{C}^8$
satisfy the commutation relation
\[
  W(v)\,W(w) = (-1)^{\omega(v,w)}\, W(w)\,W(v),
\]
so that for Lagrangian pairs ($\omega=0$) the operators commute exactly.
Combined with the Fano zero-sum property $v_1 \oplus \cdots \oplus v_4 = 0$
(which forces $W_C \in \{\pm I_8\}$) and a direct phase computation
on $|0\rangle$ (which identifies the sign as $(-1)^{\beta(C)/2}$),
this gives $W_C = s(C) I_8$ with no external citation needed.

The central circle of ideas is summarised in Figure~\ref{fig:circle}:

\begin{enumerate}
  \item \textbf{Pauli commutation} ($\omega=0$ for Lagrangian pairs)
        ensures $W_C$ is a well-defined product of commuting operators.
  \item \textbf{Fano zero-sum} forces the Pauli-string content of $W_C$
        to be $W(0) = I_8$, so $W_C \in \{\pm I_8\}$.
  \item \textbf{Phase computation} on $|0\rangle$ identifies the sign:
        $W_C|0\rangle = (-1)^{\beta(C)/2}|0\rangle$, so $W_C = s(C) I_8$.
  \item \textbf{Parity theorem} (Papers~XI--XV) gives
        $\prod_C s(C) = -1$, hence $\prod_C W_C = -I_8$.
\end{enumerate}

\noindent
The script \texttt{displacement\_operator.py} in \cite{PaperXVIsuppl}
verifies all steps computationally for all 945 contexts and 12{,}096
pentagrams.

\begin{figure}[H]
\centering
\begin{tabular}{ccccc}
Pauli commutation + Fano zero-sum
& $\xrightarrow{\text{Steps 1--2}}$
& $W_C = \pm I_8$
& $\xrightarrow{\text{Step 3 (phase on $|0\rangle$)}}$
& $W_C = s(C) \cdot I_8$ \\
& & & $\xrightarrow{\text{KS parity}}$
& $\displaystyle\prod_C W_C = -I_8$
\end{tabular}
\caption{The proof chain producing the KS obstruction
from Pauli algebra and the Fano zero-sum property.}
\label{fig:circle}
\end{figure}


\section{Weyl Operators and the Exact Commutation Relation}\label{sec:weyl}

\subsection{Weyl Operators as Pauli Strings}

Let $V = \mathbb{F}_2^6$ with the standard symplectic form
$\omega(v,w) = \sum_{i=0}^2 (v_i w_{i+3} - v_{i+3} w_i) \bmod 2$
and integer symplectic form
$\omega_{\mathrm{int}}(v,w) = \sum_{i=0}^2 (v_i w_{i+3} - v_{i+3} w_i) \in \mathbb{Z}$.

Write $v = (a_v, b_v)$ with $a_v = (v_0,v_1,v_2) \in \mathbb{F}_2^3$ (X-part)
and $b_v = (v_3,v_4,v_5) \in \mathbb{F}_2^3$ (Z-part).
The \emph{Weyl operator} $W(v)$ on $\mathbb{C}^8$
(with basis $\{|x\rangle : x \in \mathbb{F}_2^3\}$) is the
3-qubit Pauli string:
\[
  W(v) \;=\; \bigotimes_{q=0}^{2} P_q(a_q, b_q),
\]
where $P_q(0,0) = I$, $P_q(1,0) = X$, $P_q(0,1) = Z$,
$P_q(1,1) = Y = iXZ$.
Explicitly, $W(a,b)|x\rangle = i^{a \cdot b}(-1)^{b \cdot x}|x+a\rangle$
(arithmetic in $\mathbb{F}_2^3$, with $a \cdot b = \sum_i a_i b_i \in \mathbb{Z}$).

\begin{remark}[Standard Pauli operators]
The map $v \mapsto W(v)$ recovers every $n$-qubit Pauli string
(up to global phase):
$W(e_i) = X_i$, $W(e_{i+3}) = Z_i$, $W(e_i + e_{i+3}) = Y_i$,
and $W(v+w) = \pm W(v)W(w)$ modulo phase.
\end{remark}

\subsection{Pauli Commutation and the Integer Symplectic Form}

\begin{proposition}[Pauli commutation relation]\label{prop:weylcomm}
For all $v, w \in \mathbb{F}_2^6$,
\[
  W(v)\,W(w) \;=\; (-1)^{\omega(v,w)} \, W(w)\,W(v),
\]
where $\omega(v,w) \in \{0,1\}$ is the mod-$2$ symplectic form.
In particular, for Lagrangian pairs with $\omega(v,w) = 0$,
the operators $W(v)$ and $W(w)$ commute exactly.
\end{proposition}

\begin{proof}
Standard Pauli algebra: two Pauli strings commute iff they share an
even number of anticommuting single-qubit factors, which is precisely
the condition $\omega(v,w) \equiv 0 \pmod 2$.
\end{proof}

\begin{remark}[The integer symplectic form $\omega_{\mathrm{int}}$]
\label{rem:intform}
For $v, w \in L$ (Lagrangian), $\omega(v,w) = 0$ but
$\omega_{\mathrm{int}}(v,w) \in 2\mathbb{Z}$ can be $\pm 2, \pm 4, \ldots$.
The integer form records the \emph{degree of non-commutativity} in a finer
sense than the mod-2 reduction:
$W(v)W(w) = W(w)W(v)$ (exact, since $\omega = 0$),
but $W(v)W(w) \neq W(v \oplus w)$ in general when $\omega_{\mathrm{int}}(v,w) \neq 0$.

The formula $W(v)W(w) = (-i)^{\omega_{\mathrm{int}}(v,w)} W(v \oplus w)$
would be natural but is incorrect for the Pauli/Weyl operators $W(v)$
defined here (which carry the self-phase $i^{a_v \cdot b_v}$ built into
the $Y$-convention).
The correct product formula is more involved and depends on the
self-phase terms $a \cdot b$ and $a' \cdot b'$ separately
(not just their difference via $\omega_{\mathrm{int}}$).

The \emph{$\beta$-formula} $s(C) = (-1)^{\beta(C)/2}$ with
$\beta(C) = \sum_{j<k} \omega_{\mathrm{int}}(v_j, v_k)$
is not an independent postulate: it is the direct consequence of
evaluating the Pauli product $W(v_1) \cdots W(v_4)$ on the vacuum
state $|0\rangle$, where the Fano zero-sum property ensures the
final state is again $|0\rangle$ and the accumulated phase is
exactly $(-1)^{\beta(C)/2}$.
See Theorem~\ref{thm:weyl}, Step~3.
\end{remark}


\section{The Fano Zero-Sum Property}\label{sec:fano}

\begin{theorem}[Fano zero-sum]\label{thm:fano}
For every context $C = \{v_1, v_2, v_3, v_4\}$ in a Mermin pentagram,
\[
  v_1 \oplus v_2 \oplus v_3 \oplus v_4 \;=\; 0 \;\in\; \mathbb{F}_2^6.
\]
Computationally verified for all 945 distinct contexts \cite{PaperXVIsuppl}.
\end{theorem}

\begin{proof}
Each Lagrangian $L$ is a 3-dimensional isotropic subspace of
$(\mathbb{F}_2^6, \omega)$, with $|L \setminus \{0\}| = 7$ non-zero elements
forming the points of a \emph{Fano plane} (the projective plane over $\mathbb{F}_2$).
In any context $C = \{v_1, v_2, v_3, v_4\} \subset L$,
the four points are the complement of a Fano \emph{line}
$\ell = \{l_1, l_2, l_3\} \subset L$ (the unique line disjoint from $C$).

Two identities hold in $\mathbb{F}_2^6$:

\noindent
\emph{(i) XOR of all 7 non-zero Lagrangian points is $0$.}
The 7 non-zero vectors of $L \cong \mathbb{F}_2^3$ are all non-zero binary 3-tuples.
Each standard basis vector $e_i$ of $\mathbb{F}_2^3$ appears in exactly 4
of the 7 vectors (those with $e_i$-coordinate~1), so its contribution
to the total XOR is $4 \cdot e_i \equiv 0$.

\noindent
\emph{(ii) XOR of any Fano line $\{l_1, l_2, l_3\}$ is $0$.}
By the Fano collinearity condition, $l_1 + l_2 = l_3$
(mod~2 in $\mathbb{F}_2^3$), so $l_1 \oplus l_2 \oplus l_3 = 0$.

Combining: $v_1 \oplus v_2 \oplus v_3 \oplus v_4
= (l_1 \oplus l_2 \oplus l_3) \oplus (\bigoplus_{p \in L \setminus \{0\}} p)
= 0 \oplus 0 = 0$.
\end{proof}


\section{The Weyl Product Theorem}\label{sec:main}

\begin{theorem}[Weyl Product Identity]\label{thm:weyl}
For every context $C = \{v_1, v_2, v_3, v_4\}$ in a Mermin pentagram,
\[
  W_C \;:=\; W(v_1)\,W(v_2)\,W(v_3)\,W(v_4) \;=\; s(C) \cdot I_8,
\]
where $s(C) = (-1)^{\beta(C)/2} \in \{+1, -1\}$.
Computationally verified for all 945 distinct contexts \cite{PaperXVIsuppl}.
\end{theorem}

\begin{proof}
\textbf{Step 1 (Commutativity).}
Since $v_1, v_2, v_3, v_4 \in L$ (a Lagrangian), they are pairwise
isotropic: $\omega(v_j, v_k) = 0$ for all $j \neq k$.
By Proposition~\ref{prop:weylcomm}, $W(v_j)$ and $W(v_k)$ commute exactly.

\textbf{Step 2 ($W_C = \pm I_8$).}
The four operators $W(v_1), \ldots, W(v_4)$ are commuting Hermitian unitaries
each satisfying $W(v_j)^2 = I_8$ (Pauli property).
Their product $W_C$ is Hermitian (product of commuting Hermitians)
and satisfies $W_C^2 = W(v_1)^2 \cdots W(v_4)^2 = I_8$ (using commutativity
to rearrange).
Furthermore, by Theorem~\ref{thm:fano}, $v_1 \oplus \cdots \oplus v_4 = 0$.
In the Pauli group, the Pauli-string content of any product
$W(v_1) \cdots W(v_4)$ is $W(v_1 \oplus \cdots \oplus v_4) = W(0) = I_8$
up to a global phase; hence $W_C \in \mathbb{C} \cdot I_8$.
Combined with $W_C = \pm I_8$ (from Step~1 above), we obtain $W_C \in \{+I_8, -I_8\}$.

\textbf{Step 3 (Identification of the sign).}
Write $v_j = (a_j, b_j)$ with $a_j, b_j \in \mathbb{F}_2^3$.
Using $W(a,b)|x\rangle = i^{a \cdot b}(-1)^{b \cdot x}|x \oplus a\rangle$
(all dot products integer-valued), we compute $W_C|0\rangle$
by applying the four operators from right to left.
The Fano zero-sum property (Theorem~\ref{thm:fano}) gives
$a_1 \oplus a_2 \oplus a_3 \oplus a_4 = 0$, so the final state
is again $|0\rangle$ and
\[
  W_C |0\rangle \;=\; c \cdot |0\rangle,
\]
where the accumulated phase is
\[
  c \;=\; i^{\sum_j a_j \cdot b_j} \cdot (-1)^{\sum_{j<k} b_j \cdot a_k}.
\]
(Since all entries of $a_j$, $b_j$ are binary, correction terms
arising from XOR versus integer addition in the exponents of $(-1)$
are always even integers and do not contribute to the phase.)
Let $S = \sum_j a_j \cdot b_j$, $T = \sum_{j<k} b_j \cdot a_k$,
and $U = \sum_{j<k} a_j \cdot b_k$.
The key identity (integer arithmetic) is
\[
  S + U + T \;=\; \Bigl(\sum_j a_j\Bigr) \cdot \Bigl(\sum_k b_k\Bigr).
\]
The Fano zero-sum $v_1 \oplus \cdots \oplus v_4 = 0$ in $\mathbb{F}_2^6$
implies that each component of $\sum_j a_j$ and $\sum_k b_k$ is an
even integer (each is $0$, $2$, or $4$), so their dot product is
divisible by~4.
Hence $S + U + T \equiv 0 \pmod 4$, giving $i^S = i^{-(U+T)}$.
The phase simplifies to
\[
  c \;=\; i^{-(U+T)} \cdot (-1)^T
  \;=\; i^{-(U+T)} \cdot i^{2T}
  \;=\; i^{T-U}
  \;=\; i^{-\beta(C)},
\]
where $\beta(C) = U - T = \sum_{j<k} \omega_{\mathrm{int}}(v_j, v_k)$.
Since $\beta(C) \in 2\mathbb{Z}$ (Remark~\ref{rem:betaeven}),
\[
  c \;=\; i^{-\beta(C)} \;=\; (-1)^{\beta(C)/2} \;=\; s(C).
\]
Since $W_C = c \cdot I_8$ (Step~2), we obtain $W_C = s(C) \cdot I_8$.
\end{proof}

\begin{remark}[Why $\beta(C)$ is always even]\label{rem:betaeven}
For $v_j, v_k \in L$, each term $\omega_{\mathrm{int}}(v_j,v_k)$ is
individually even: since $\omega = \omega_{\mathrm{int}} \bmod 2$,
the isotropic condition $\omega(v_j,v_k) = 0$ forces
$\omega_{\mathrm{int}}(v_j,v_k) \equiv 0 \pmod 2$.
Hence $\beta(C) = \sum_{j<k} \omega_{\mathrm{int}}(v_j,v_k)$ is a
sum of even integers, so $\beta(C) \in 2\mathbb{Z}$ and
$s(C) = (-1)^{\beta(C)/2} \in \{+1,-1\}$.
\end{remark}

\begin{remark}[Context-level vs.\ Lagrangian-level identification]
Paper~XV showed that the Lagrangian-level conjecture
$\beta(C)/2 \equiv c(g_C, g_C^{-1}) \pmod 2$ fails because 60\% of
Lagrangians have varying $h$-values across their contexts \cite{PaperXV}.
Theorem~\ref{thm:weyl} explains why: the correct encoding of $\beta(C)$
is \emph{context-level} because $W_C$ depends on the four specific rays
$v_1, \ldots, v_4$, not just on the Lagrangian $L_C$ that contains them.
Different contexts in the same Lagrangian generally have different
$\beta$-values, reflected in different $W_C$ signs.
\end{remark}


\section{The KS Obstruction as a Weyl-Algebra Identity}\label{sec:ks}

\begin{corollary}[Pentagram product identity]\label{cor:pent}
For every Mermin pentagram $\mathbf{P} = (C_1, \ldots, C_5)$,
\[
  \prod_{i=1}^{5} W_{C_i} \;=\; -I_8.
\]
Computationally verified for all 12{,}096 Mermin pentagrams
\cite{PaperXVIsuppl}.
\end{corollary}

\begin{proof}
By Theorem~\ref{thm:weyl}, $W_{C_i} = s(C_i) I_8$ for each $i$.
The product of five scalar multiples of $I_8$ is
$\bigl(\prod_{i=1}^5 s(C_i)\bigr) I_8$.
The parity theorem \cite{PaperXI} gives $\prod_{i=1}^5 s(C_i) = -1$
for every Mermin pentagram.
\end{proof}

\begin{remark}[The mod~4 structure from the $\beta$-formula]
The equation $\prod_C W_C = -I_8$ is equivalent to $\beta_{\mathrm{sum}} \equiv 2 \pmod 4$:
the product of $5$ context signs is $(-1)^{\beta_{\mathrm{sum}}/2}$,
which equals $-1$ iff $\beta_{\mathrm{sum}}/2$ is odd, i.e., $\beta_{\mathrm{sum}} \equiv 2 \pmod 4$.
The derivation above makes the origin of the mod~4 structure explicit:
it arises from the $\beta$-formula $s(C) = (-1)^{\beta(C)/2}$ where
$\beta(C) = \sum_{j<k} \omega_{\mathrm{int}}(v_j, v_k)$ uses the
\emph{integer} symplectic form, not its mod-2 reduction.
The mod-2 reduction loses the phase information that encodes the sign.
\end{remark}


\section{G\texorpdfstring{$_2(2)$}{2(2)} Equivariance and the Orbit Structure}\label{sec:g2}

The Weyl Product Identity is compatible with the $G_2(2)$ orbit structure
established in Papers~X--XV.

\begin{proposition}[$G_2(2)$-equivariance of Weyl signs]\label{prop:equiv}
For any $g \in G_2(2) \leq Sp(6,\mathbb{F}_2)$,
$s(g \cdot C) = s(C)$: the context sign is a $G_2(2)$-invariant.
Hence $W_{g \cdot C} = s(C) \cdot I_8 = W_C$.
\end{proposition}

\begin{proof}
$G_2(2)$ preserves the symplectic form, hence $\beta(g \cdot C) = \beta(C)$
(since $\omega_{\mathrm{int}}(gv, gw) = \omega_{\mathrm{int}}(v,w)$ for
$g \in Sp(6,\mathbb{F}_2)$), and therefore $s(g \cdot C) = s(C)$.
\end{proof}

\begin{table}[H]
\centering
\caption{Weyl sign distribution for the 945 contexts, broken down by
$k$-type $k = |L_C \cap V|$ of the containing Lagrangian.
$s(C) = W_C[0,0]$ as computed by \texttt{displacement\_operator.py}
\cite{PaperXVIsuppl}.}
\label{tab:weylsigns}
\begin{tabular}{lrrrr}
\toprule
$k$-type & Total & $s(C) = -1$ & $s(C) = +1$ & \% minus \\
\midrule
$k = 0$ & 448 & 176 & 272 & 39.3\% \\
$k = 1$ & 392 & 128 & 264 & 32.7\% \\
$k = 3$ &  98 &  20 &  78 & 20.4\% \\
$k = 7$ &   7 &   0 &   7 &  0.0\% \\
\midrule
Total   & 945 & 324 & 621 & 34.3\% \\
\bottomrule
\end{tabular}
\end{table}

\begin{remark}[$k = 7$ and the uniform zero case]
For $k = 7$, the Lagrangian $L_C = V$ itself.
Every context $C \subset V$ has $\omega_{\mathrm{int}}(v,w) = 0$ for
$v, w \in V = \mathbb{F}_2^3 \times \{0\}^3$ (since $b_v = b_w = 0$),
so $\beta(C) = 0$ and $s(C) = +1$ uniformly.
The 0\% minus rate for $k=7$ in Table~\ref{tab:weylsigns} is an algebraic
consequence, not an empirical coincidence.
\end{remark}


\section{Resolution of Previous Open Problems}\label{sec:resolution}

Theorem~\ref{thm:weyl} and Corollary~\ref{cor:pent} complete the
research arc of Papers~XI--XV.
Table~\ref{tab:resolution} summarises the resolution.

\begin{table}[H]
\centering
\caption{Resolution of open problems from Papers~XI--XV by the Weyl
Product Identity (Theorem~\ref{thm:weyl}).}
\label{tab:resolution}
\begin{tabular}{p{1.4cm}p{5.2cm}p{5.2cm}}
\toprule
Paper & Open Problem & Paper~XVI Resolution \\
\midrule
XI & Origin of the $\beta$-formula
    & $\beta(C) = \sum_{j<k}\omega_{\mathrm{int}}(v_j,v_k)$ is the
      phase accumulated by $W_C$ acting on $|0\rangle$;
      $W_C = s(C)\cdot I_8$ derived algebraically (no external citation) \\
\addlinespace
XII & Origin of mod~4 structure
    & The $\beta$-formula uses $\omega_{\mathrm{int}} \in \mathbb{Z}$
      (not mod~2); the sign $(-1)^{\beta/2}$ is intrinsically mod~4 \\
\addlinespace
XIII & Why Maslov index fails
    & Maslov index is a mod-2 construction; it cannot encode the
      integer symplectic phase $\omega_{\mathrm{int}}$ that determines
      the sign via the $\beta$-formula \\
\addlinespace
XIV & $h$-pattern structure
    & $h(C) = \beta(C)/2 \bmod 2 = $ sign of $W_C$;
      the parity structure reflects the Pauli algebra of the context \\
\addlinespace
XV  & Context-level $\beta$ identification
     (Open Problem~1)
    & \textbf{Solved:} $\beta(C)/2 \bmod 2 = $ sign of $W_C$;
      the Pauli product $W_C$ is the context-level object \\
\addlinespace
XV  & Why $\beta$-cocycle conjecture failed
    & $\beta(C)$ is determined by the Pauli product sign,
      not a metaplectic cocycle; cocycles are Lagrangian-level,
      Pauli signs are context-level \\
\bottomrule
\end{tabular}
\end{table}


\section{Open Problems}\label{sec:open}

\begin{enumerate}
  \item \textbf{$n$-qubit generalization}:
        Does the Weyl Product Identity extend to Mermin pentagrams in
        $W(2n-1, \mathbb{F}_2)$ for arbitrary $n$?
        The Fano zero-sum and Weyl commutation both generalise to
        arbitrary Lagrangians, but the context structure and parity theorem
        need to be established for general $n$.

  \item \textbf{Gerbe / Dixmier--Douady class}:
        The sign assignment $C \mapsto s(C) \in \{+1,-1\}$
        with $\prod_{C \in \mathbf{P}} s(C) = -1$ is a
        \v{C}ech 1-cocycle on the pentagram nerve.
        Is this cocycle the Dixmier--Douady class of a gerbe on the
        Lagrangian Grassmannian
        \cite{DixmierDouady,Brylinski1993}?

  \item \textbf{Physical interpretation of $W_C = s(C)I_8$}:
        The identity $W(v_1) \cdots W(v_4) = \pm I$ is the
        stabilizer-group statement that the product of all generators
        of a maximal abelian subgroup is $\pm I$.
        Does the specific sign $s(C) = (-1)^{\beta(C)/2}$ have a
        physical interpretation in terms of measurement backaction or
        contextuality witnesses?

  \item \textbf{Algebraic proof of $\prod_C s(C) = -1$ from Weyl alone
        (primary open problem)}:
        Corollary~\ref{cor:pent} derives $\prod_C W_C = -I_8$
        by invoking the parity theorem of Paper~XI \cite{PaperXI}
        (which is itself a computational result over all 12{,}096 pentagrams).
        Theorem~\ref{thm:weyl} is now algebraically complete at the
        \emph{context} level: both $W_C = \pm I_8$ and
        $\mathrm{sign}(W_C) = s(C)$ are derived from Pauli algebra
        and the Fano zero-sum property without external citation.
        The remaining gap is a direct Weyl-algebraic proof that
        $\prod_{C \in \mathbf{P}} W_C = -I_8$ without invoking the
        12{,}096-pentagram enumeration.
        This is the sole remaining obstacle to full algebraic
        completeness of the mod~4 problem.
        The $G_2(2)$ orbit structure (Papers~X--XV) and the
        representation-theoretic properties of the Weil representation
        are the most promising entry points.

  \item \textbf{Orbit~4 anomaly via Weyl signs}:
        Paper~XV showed that $O_4$ has a 60.6\% 1-minus rate vs.\
        65--68\% for $O_1$--$O_3$, with the asymmetric metaplectic
        lift as the structural candidate.
        Does the Weyl sign distribution of Table~\ref{tab:weylsigns},
        broken down by orbit, show a differential pattern for $O_4$
        that can be derived from the NS-asymmetric $S_3$ lift?
\end{enumerate}


\appendix
\setcounter{table}{4}
\renewcommand{\thetable}{A\arabic{table}}
\renewcommand{\theHtable}{A\arabic{table}}
\section{Rigor Self-Assessment}

\begin{table}[H]
\centering
\caption{Rigor classification and verification scope for all main results.}
\label{tab:rigor}
\begin{tabular}{p{6cm}lp{4cm}}
\toprule
Result & Type & Scope \\
\midrule
Prop~\ref{prop:weylcomm} (Pauli commutation $\omega(v,w)=0$)
  & Algebraic & Standard Pauli algebra \\
Rem~\ref{rem:intform} ($\omega_{\mathrm{int}}$ vs $\omega$ mod~2)
  & Algebraic & Structural note; no computation \\
Thm~\ref{thm:fano} (Fano zero-sum)
  & Algebraic & Exact proof + all 945 contexts \cite{PaperXVIsuppl} \\
Rem~\ref{rem:betaeven} ($\beta(C)$ always even)
  & Algebraic & Follows from Prop~\ref{prop:weylcomm} \\
Thm~\ref{thm:weyl} (Weyl Product Identity)
  & Algebraic & Exact proof (Steps 1--3); \\
& & verified 945/945 \cite{PaperXVIsuppl} \\
Cor~\ref{cor:pent} (Pentagram product $= -I_8$)
  & Algebraic + Cite & Follows from Thm~\ref{thm:weyl} + \cite{PaperXI}; \\
& & 12,096 pentagrams verified \cite{PaperXVIsuppl} \\
Prop~\ref{prop:equiv} ($G_2(2)$-equivariance)
  & Algebraic & Follows from Sp-invariance of $\omega_{\mathrm{int}}$ \\
Table~\ref{tab:weylsigns} (Weyl sign by $k$-type)
  & Computational & All 945 contexts \cite{PaperXVIsuppl} \\
\bottomrule
\end{tabular}
\end{table}


\begin{thebibliography}{99}

\bibitem[PaperX]{PaperX}
Z.-L. Chen,
``The Equiangular Characterization of Mermin Pentagrams:
Symplectic Embedding, 10-Ray Cap, and the Failure of the
Collinearity--Obstruction Dictionary,''
\textit{Zenodo, DOI: 10.5281/zenodo.20476659} (2026).

\bibitem[PaperXI]{PaperXI}
Z.-L. Chen,
``Quadratic Refinement and the Parity of Mermin Pentagrams:
The $\beta$-Formula, Geometric Decomposition, and the
Kochen--Specker Obstruction from Equiangular Geometry,''
\textit{Zenodo, DOI: 10.5281/zenodo.20482595} (2026).

\bibitem[PaperXII]{PaperXII}
Z.-L. Chen,
``The T-Vector Theorem: Symplectic Characteristic Elements
and the Kochen--Specker Obstruction,''
\textit{Zenodo, DOI: 10.5281/zenodo.20490118} (2026).

\bibitem[PaperXIII]{PaperXIII}
Z.-L. Chen,
``The Maslov Index and the Kochen--Specker Obstruction:
Negative Results, the $k$-Profile Theorem, and Orbit~4 Anomaly,''
\textit{Zenodo, DOI: 10.5281/zenodo.20496513} (2026).

\bibitem[PaperXIV]{PaperXIV}
Z.-L. Chen,
``Stabilizer Algebra and the $k$-Profile Theorem for Mermin Pentagrams,''
\textit{Zenodo, DOI: 10.5281/zenodo.20502868} (2026).

\bibitem[PaperXV]{PaperXV}
Z.-L. Chen,
``The Weil Representation and the $S_3$ Lifting Classification
for Mermin Pentagrams,''
\textit{Zenodo, DOI: 10.5281/zenodo.20519654} (2026).

\bibitem[PaperXVIsuppl]{PaperXVIsuppl}
Z.-L. Chen,
``Supplementary computational scripts for Paper~XVI,''
\textit{Zenodo, DOI: 10.5281/zenodo.20519672} (2026).

\bibitem[DixmierDouady]{DixmierDouady}
J. Dixmier and A. Douady,
``Champs continus d'espaces hilbertiens et de $C^*$-alg\`ebres,''
\textit{Bull. Soc. Math. France} \textbf{91} (1963), 227--284.

\bibitem[Brylinski1993]{Brylinski1993}
J.-L. Brylinski,
\textit{Loop Spaces, Characteristic Classes and Geometric Quantization},
Birkh\"auser, Boston (1993).

\bibitem[Gerardin]{Gerardin}
P. G\'erardin,
``Weil representations associated to finite fields,''
\textit{J. Algebra} \textbf{46} (1977), 54--101.

\bibitem[Mermin]{Mermin}
N.D. Mermin,
``Hidden variables and the two theorems of John Bell,''
\textit{Rev. Mod. Phys.} \textbf{65} (1993), 803--815.

\end{thebibliography}

\end{document}
